% ========================================
% NOMA-GNN Paper Figures - INTEGRATED INTO main.tex
% All 5 TikZ figures have been added to the IEEE paper
% ========================================

% Figure locations in main.tex:
% 1. graph_construction.tex - After "Problem Reformulation as Graph Learning" (Section III.A)
% 2. gnn_architecture.tex - After describing multi-task decoder heads (Section III.B)
% 3. graphsage_detail.tex - After explaining MEAN aggregator (Section III.B.1)
% 4. inference_pipeline.tex - After inference complexity discussion (Section III.E)
% 5. training_inference_comparison.tex - After training details (Section III.D)

% ========================================
% NOMA-GNN Paper Figures - Usage Guide
% Place these in your LaTeX document where appropriate

% ========================================
% Figure 1: GNN Architecture Overview
% Best location: Section III (GNN Methodology)
% ========================================
\begin{figure}[t]
\centering
% GNN Architecture Diagram - Compact Version
\begin{tikzpicture}[
    node/.style={circle, draw=black!90, fill=blue!35, very thick, minimum size=9mm, font=\small},
    feature/.style={rectangle, draw=black!80, fill=green!25, very thick, minimum width=18mm, minimum height=7mm, font=\small, align=center},
    layer/.style={rectangle, draw=black!80, fill=orange!25, very thick, minimum width=26mm, minimum height=9mm, font=\small\bfseries, align=center},
    decoder/.style={rectangle, draw=black!80, fill=purple!25, very thick, minimum width=22mm, minimum height=8mm, font=\small, align=center},
    arrow/.style={->, >=stealth, very thick},
    scale=0.88, every node/.style={scale=0.88}
]

% Input Layer - Graph Structure (Left Side)
\node[font=\normalsize\bfseries] at (0, 6) {Input Graph};
\node[node] (u1) at (-0.7, 5) {$u_1$};
\node[node] (u2) at (0.7, 5) {$u_2$};
\node[node] (u3) at (-0.7, 3.8) {$u_3$};
\node[node] (u4) at (0.7, 3.8) {$u_4$};

% Edges
\draw[very thick, gray!70] (u1) -- (u2);
\draw[very thick, gray!70] (u1) -- (u3);
\draw[very thick, gray!70] (u2) -- (u4);
\draw[very thick, gray!70] (u3) -- (u4);

% Node Features
\node[font=\small, below=10mm of u3, text width=30mm, align=center] {Features: $|h_i|^2, d_i, \theta_i$};

% GraphSAGE Encoder (Center)
\begin{scope}[xshift=5.5cm]
    \node[font=\normalsize\bfseries] at (0, 6) {Encoder};
    \node[layer] (sage1) at (0, 5.2) {GraphSAGE\\Layer 1};
    \node[layer, below=7mm of sage1] (sage2) {GraphSAGE\\Layer 2};
    \node[layer, below=7mm of sage2] (sage3) {GraphSAGE\\Layer 3};
    
    \draw[arrow] (sage1) -- node[right, font=\scriptsize] {$\mathbf{h}^{(1)}$} (sage2);
    \draw[arrow] (sage2) -- node[right, font=\scriptsize] {$\mathbf{h}^{(2)}$} (sage3);
    
    \node[feature, below=3mm of sage3, fill=cyan!30] (emb) {Embeddings\\$\mathbb{R}^{128}$};
    \draw[arrow] (sage3) -- (emb);
\end{scope}

% Connection from input to encoder
\draw[arrow, dashed, gray!60] (u2) -- (sage1);
\draw[arrow, dashed, gray!60] (u3) -- (sage1);

% Edge Processing (Right Center)
\begin{scope}[xshift=10cm]
    \node[feature, fill=yellow!30, minimum width=24mm] (concat) at (0, 3.3) {Edge Concat\\$[\mathbf{h}_i \| \mathbf{h}_j]$};
\end{scope}

\draw[arrow] (emb) -| (concat);

% Multi-Task Decoders (Right Side)
\begin{scope}[xshift=14cm]
    \node[font=\normalsize\bfseries] at (0, 6) {Decoders};
    
    \node[decoder] (pair) at (0, 5) {Pairing\\Classifier};
    \node[decoder, below=7mm of pair] (rate) {Sum-Rate\\Regressor};
    \node[decoder, below=7mm of rate] (power) {Power\\Regressor};
    
    % Outputs
    \node[feature, right=9mm of pair, fill=red!30, minimum width=16mm] (out1) {$p_{ij}$};
    \node[feature, right=9mm of rate, fill=red!30, minimum width=16mm] (out2) {$\hat{r}_{ij}$};
    \node[feature, right=9mm of power, fill=red!30, minimum width=16mm] (out3) {$\hat{\alpha}_{ij}$};
    
    \draw[arrow] (pair) -- (out1);
    \draw[arrow] (rate) -- (out2);
    \draw[arrow] (power) -- (out3);
    
    % Output labels
    \node[font=\small, right=2mm of out1] {Pairing};
    \node[font=\small, right=2mm of out2] {Rate};
    \node[font=\small, right=2mm of out3] {Power};
\end{scope}

% Connections from concat to decoders
\draw[arrow] (concat.east) -- ++(0.5,0) |- (pair.west);
\draw[arrow] (concat.east) -- (rate.west);
\draw[arrow] (concat.east) -- ++(0.5,0) |- (power.west);

% Complexity annotation at bottom
\node[draw, very thick, fill=gray!15, rectangle, text width=100mm, font=\small, align=center] at (7.5, -.8) {
\textbf{Complexity:} Training $O(E \cdot D^2)$ — Inference $O(N \log N)$
};

\end{tikzpicture}

\caption{NOMA-GNN architecture: GraphSAGE encoder with multi-task decoders. Node features include channel gain $|h_i|^2$, distance $d_i$, angular position $\theta_i$, and SNR. Three GraphSAGE layers produce 128-dimensional embeddings, which are concatenated pairwise and fed to three task-specific MLPs for pairing classification, sum-rate regression, and power allocation.}
\label{fig:gnn_architecture}
\end{figure}

% ========================================
% Figure 2: Graph Construction Process
% Best location: Section III.A (Problem Reformulation)
% ========================================
\begin{figure}[t]
\centering
% Graph Construction Process
\begin{tikzpicture}[
    user/.style={circle, draw=black!80, fill=blue!30, thick, minimum size=10mm, font=\footnotesize},
    bs/.style={regular polygon, regular polygon sides=3, draw=black!80, fill=red!50, thick, minimum size=12mm},
    edge/.style={thick, draw=green!60!black},
    noedge/.style={thick, draw=red!60, dashed},
    arrow/.style={->, >=stealth, thick},
    scale=0.9, every node/.style={scale=0.9}
]

% Base Station
\node[bs] (bs) at (0, 0) {};
\node[font=\small\bfseries, below=1mm of bs] {BS};

% Users
\node[user] (u1) at (-2.5, 2) {$u_1$};
\node[user] (u2) at (2.5, 2.5) {$u_2$};
\node[user] (u3) at (-1.5, -2) {$u_3$};
\node[user] (u4) at (2, -2.2) {$u_4$};

% Channel info
\node[font=\tiny, above=3mm of u1] {Strong};
\node[font=\tiny, text=green!50!black] at ([yshift=3mm]u1.north) {$h_1 = -62$ dB};

\node[font=\tiny, above=3mm of u2] {Strong};
\node[font=\tiny, text=green!50!black] at ([yshift=3mm]u2.north) {$h_2 = -58$ dB};

\node[font=\tiny, below=1mm of u3] {Weak};
\node[font=\tiny, text=red!60!black] at ([yshift=-5mm]u3.south) {$h_3 = -98$ dB};

\node[font=\tiny, below=1mm of u4] {Medium};
\node[font=\tiny, text=orange!80!black] at ([yshift=-5mm]u4.south) {$h_4 = -75$ dB};

% Lines to BS
\draw[gray, thin, dashed] (bs) -- (u1);
\draw[gray, thin, dashed] (bs) -- (u2);
\draw[gray, thin, dashed] (bs) -- (u3);
\draw[gray, thin, dashed] (bs) -- (u4);

% Valid edges (SIC feasible + angular separation)
\draw[edge, line width=1.5pt] (u1) to[bend left=15] node[midway, above, font=\tiny, fill=white] {Valid} (u3);
\draw[edge, line width=1.5pt] (u2) to[bend right=15] node[midway, right, font=\tiny, fill=white] {Valid} (u4);

% Invalid edge (too close in channel gain)
\draw[noedge, line width=1.2pt] (u1) to[bend right=10] node[midway, above, font=\tiny, fill=white] {$\times$ Similar $h$} (u2);

% Constraint boxes
\node[draw, rectangle, thick, fill=yellow!20, text width=32mm, font=\tiny, align=left] at (-4.5, 0) {
\textbf{Edge Constraints:}\\
$\bullet$ $|h_i|^2 / |h_j|^2 \geq \gamma_{\text{SIC}}$\\
$\bullet$ $|\theta_i - \theta_j| \geq 25^\circ$\\
$\bullet$ $|h_i|^2 < |h_j|^2$
};

% Graph representation
\begin{scope}[xshift=8cm, yshift=0.5cm]
    \node[font=\small\bfseries] at (0, 2.8) {Graph $\mathcal{G}$};
    
    % Nodes
    \node[user, minimum size=8mm] (g1) at (-1, 1.5) {$v_1$};
    \node[user, minimum size=8mm] (g2) at (1, 1.5) {$v_2$};
    \node[user, minimum size=8mm] (g3) at (-1, -0.5) {$v_3$};
    \node[user, minimum size=8mm] (g4) at (1, -0.5) {$v_4$};
    
    % Edges with labels
    \draw[edge, line width=1.5pt] (g1) -- (g3) node[midway, left, font=\tiny] {$e_{13}$};
    \draw[edge, line width=1.5pt] (g2) -- (g4) node[midway, right, font=\tiny] {$e_{24}$};
    
    % Features
    \node[font=\tiny, align=center, below=8mm of g3] {
        $\mathbf{X} \in \mathbb{R}^{N \times 4}$\\
        $\mathcal{E}$ = feasible pairs
    };
\end{scope}

% Arrow indicating transformation
\draw[->, ultra thick, gray] (4.5, 0) -- (6.5, 0) node[midway, above, font=\small] {Construct};

\end{tikzpicture}

\caption{Graph construction from channel state information. Users with feasible NOMA pairing (satisfying SIC threshold, angular separation, and channel gain ordering) are connected by edges. The resulting graph $\mathcal{G}(\mathcal{V}, \mathcal{E})$ captures physical constraints and serves as input to the GNN.}
\label{fig:graph_construction}
\end{figure}

% ========================================
% Figure 3: Inference Pipeline
% Best location: Section III.E (Inference and Greedy Matching)
% ========================================
\begin{figure}[t]
\centering
% Inference Pipeline Flow
\begin{tikzpicture}[
    box/.style={rectangle, draw=black!80, thick, minimum width=28mm, minimum height=10mm, font=\footnotesize, align=center},
    data/.style={box, fill=blue!20},
    process/.style={box, fill=orange!20},
    output/.style={box, fill=green!20},
    decision/.style={diamond, draw=black!80, thick, fill=yellow!20, minimum size=12mm, font=\tiny, align=center, aspect=2},
    arrow/.style={->, >=stealth, thick},
    scale=0.88, every node/.style={scale=0.88}
]

% Column 1: Input
\node[data] (csi) at (0, 0) {Channel State\\Information\\$\{h_i, d_i, \theta_i\}$};

\node[process, below=10mm of csi] (graph) {Graph\\Construction\\$\mathcal{G}(\mathcal{V}, \mathcal{E})$};

\draw[arrow] (csi) -- (graph);

% Column 2: GNN Processing
\node[process, below=10mm of graph] (gnn) {GNN Forward\\Pass\\(GraphSAGE)};

\draw[arrow] (graph) -- node[right, font=\tiny] {$\mathbf{X}, \mathcal{E}$} (gnn);

\node[data, below=10mm of gnn] (embed) {Node\\Embeddings\\$\{\mathbf{h}_i\}_{i=1}^N$};

\draw[arrow] (gnn) -- (embed);

% Column 3: Scoring
\node[process, below=10mm of embed] (score) {Edge Scoring\\$\forall (i,j) \in \mathcal{E}$};

\draw[arrow] (embed) -- (score);

\node[data, below=10mm of score] (scores) {Pairing Scores\\$\{p_{ij}, \hat{r}_{ij}, \hat{\alpha}_{ij}\}$};

\draw[arrow] (score) -- (scores);

% Column 4: Filtering
\node[decision, below=15mm of scores] (filter) {SIC\\Valid?};

\draw[arrow] (scores) -- (filter);

\node[data, left=8mm of filter] (discard) {Discard};
\node[process, right=8mm of filter] (keep) {Candidate\\Set $\mathcal{E}'$};

\draw[arrow] (filter) -- node[above, font=\tiny] {No} (discard);
\draw[arrow] (filter) -- node[above, font=\tiny] {Yes} (keep);

% Column 5: Matching
\node[process, below=12mm of filter] (greedy) {Greedy\\Matching\\(Sort by $p_{ij}$)};

\draw[arrow] (keep) |- (greedy);

\node[output, below=10mm of greedy] (pairs) {Final Pairs\\$\mathcal{M}^*$};

\draw[arrow] (greedy) -- (pairs);

% Timing annotations
\node[font=\tiny, gray, right=2mm of graph] {$O(N^2)$};
\node[font=\tiny, gray, right=2mm of gnn] {$O(E \cdot D)$};
\node[font=\tiny, gray, right= 20mm of greedy] {$O(E \log E)$};

% Overall complexity
\node[draw, thick, fill=cyan!10, rectangle, text width=28mm, font=\tiny, align=center] at (4, -8) {
\textbf{Total Complexity:}\\
$O(N \log N)$\\[1mm]
\textbf{Runtime (N=500):}\\
752 ms
};

% Side annotations for clarity
\node[font=\footnotesize\bfseries, left=15mm of csi, rotate=90, anchor=south] {Input};
\node[font=\footnotesize\bfseries, left=15mm of gnn, rotate=90, anchor=south] {Encoding};
\node[font=\footnotesize\bfseries, left=15mm of score, rotate=90, anchor=south] {Scoring};
\node[font=\footnotesize\bfseries, left=15mm of greedy, rotate=90, anchor=south] {Matching};
\node[font=\footnotesize\bfseries, left=15mm of pairs, rotate=90, anchor=south] {Output};

\end{tikzpicture}

\caption{Real-time inference pipeline. Given new CSI, the system constructs a graph, performs GNN forward pass, scores candidate pairs, filters by SIC constraints, and executes greedy matching. Total complexity is $O(N \log N)$ with 752 ms runtime for N=500 users.}
\label{fig:inference_pipeline}
\end{figure}

% ========================================
% Figure 4: GraphSAGE Message Passing
% Best location: Section III.B (GraphSAGE Encoder)
% ========================================
\begin{figure}[t]
\centering
% GraphSAGE Message Passing Illustration
\begin{tikzpicture}[
    >=stealth,
    arrow/.style={->,thick},
    message/.style={->,thick,dashed,orange!80},
    node/.style={circle,draw=black!80,fill=blue!25,thick,minimum size=12mm,font=\footnotesize},
    neighbor/.style={circle,draw=black!80,fill=gray!15,thick,minimum size=11mm,font=\footnotesize},
    feature/.style={rectangle,draw=black!70,fill=green!15,minimum width=19mm,minimum height=7mm,font=\scriptsize},
    every node/.style={transform shape},
    scale=0.95
]

%--------------------------------------------------------------------
% Left: layer \ell=0, local neighborhood
%--------------------------------------------------------------------
\node[font=\large\bfseries] at (-2.8,3.2) {Layer $\ell = 0$};

\node[node]     (u0) at (0,1.8) {$u_i$};
\node[neighbor] (n1) at (-2,0.5) {$u_1$};
\node[neighbor] (n2) at (0,0.1)   {$u_2$};
\node[neighbor] (n3) at ( 2,0.5) {$u_3$};

\node[feature,below=18mm of u0] (f0) {$\mathbf{h}_i^{(0)}$};
\node[feature,below=3mm of n1] (f1) {$\mathbf{h}_1^{(0)}$};
\node[feature,below=11mm of n2] (f2) {$\mathbf{h}_2^{(0)}$};
\node[feature,below=1.5mm of n3] (f3) {$\mathbf{h}_3^{(0)}$};

% Edges in the original graph
\draw[thick,gray] (u0) -- (n1);
\draw[thick,gray] (u0) -- (n2);
\draw[thick,gray] (u0) -- (n3);

\node[font=\scriptsize,gray] at (1.4,2.4) {$\mathcal{N}(u_i)$};

%--------------------------------------------------------------------
% Middle: aggregation block
%--------------------------------------------------------------------
\begin{scope}[xshift=6.2cm]
  \node[font=\large\bfseries] at (0,3.2) {Aggregation};

  \node[draw,thick,fill=yellow!15,
        rectangle,
        minimum width=40mm,
        minimum height=15mm,
        font=\scriptsize,
        align=center,
        text width=42mm] (agg) at (0,1.2) {%
        \textbf{MEAN}\\[1mm]
        $\displaystyle \mathbf{m}_i^{(\ell)}
        = \frac{1}{|\mathcal{N}(i)|}
          \sum_{j\in\mathcal{N}(i)} \mathbf{h}_j^{(\ell-1)}$};

  % incoming neighbor messages
  \draw[message] (f1.east) to[bend left=18] (agg.west);
  \draw[message] (f2.east) to[bend left=6]  (agg.west);
  \draw[message] (f3.east) to[bend right=18] (agg.south west);

  \node[font=\scriptsize,text=orange!80] at (0,2.4) {Neighbor messages};
\end{scope}

%--------------------------------------------------------------------
% Right: update block
%--------------------------------------------------------------------
\begin{scope}[xshift=12.2cm]
  \node[font=\large\bfseries] at (0,3.2) {Update};

  \node[draw,thick,fill=cyan!18,
        rectangle,
        minimum width=42mm,
        minimum height=19mm,
        font=\scriptsize,
        align=center,
        text width=44mm] (update) at (0,1.0) {%
        \textbf{Concatenate:}\\[1mm]
        $[\mathbf{h}_i^{(\ell-1)} \,\|\, \mathbf{m}_i^{(\ell)}]$\\[2mm]
        \textbf{Transform:}\\[1mm]
        $\mathbf{h}_i^{(\ell)}
          = \sigma\!\bigl(\mathbf{W}^{(\ell)} [\cdot]
          + \mathbf{b}^{(\ell)}\bigr)$};

  % self feature and aggregated message to update
  \draw[arrow] (f0.east) to[bend right=25]
      node[midway,above,font=\scriptsize] {self} (update.west);
  \draw[arrow] (agg.east) -- (update.west);
\end{scope}

%--------------------------------------------------------------------
% Output embedding
%--------------------------------------------------------------------
\begin{scope}[xshift=12.2cm,yshift=-3.1cm]
  \node[feature,fill=purple!25,minimum width=22mm,minimum height=8mm,font=\scriptsize]
        (out) at (0,0) {$\mathbf{h}_i^{(\ell)}$};
  \draw[arrow,ultra thick] (update.south) -- (out.north);

  \node[font=\scriptsize,below=2mm of out] {Updated embedding};
\end{scope}

%--------------------------------------------------------------------
% Horizontal process-flow labels
%--------------------------------------------------------------------
\draw[->,ultra thick,gray!50]
      (-2.8,-2.8) -- (4.2,-2.8)
      node[midway,below=1mm,font=\small] {Message Passing};

\draw[->,ultra thick,gray!50]
      (6.0,-1.6) -- (11.5,-1.6)
      node[midway,below=1mm,font=\small] {Feature Update};

%--------------------------------------------------------------------
% Iteration note box
%--------------------------------------------------------------------
\node[draw,thick,fill=orange!8,
      rectangle,
      text width=60mm,
      font=\scriptsize,
      align=left] at (5.8,-3.7) {%
  \textbf{Repeat for $L = 3$ layers:}\\
  $\bullet$ Each layer aggregates from 1-hop neighbors\\
  $\bullet$ Layer 3 captures 3-hop neighborhood\\
  $\bullet$ Hidden dimension: 128
};

\end{tikzpicture}

\caption{GraphSAGE message passing mechanism. At each layer $\ell$, node $u_i$ aggregates features from neighbors $\mathcal{N}(i)$ using MEAN aggregation, concatenates with its own features, and applies a learnable transformation. Three layers enable each node to incorporate 3-hop neighborhood information.}
\label{fig:graphsage_detail}
\end{figure}

% ========================================
% Figure 5: Training vs Inference Comparison
% Best location: Section IV (Results) or Section III.E
% ========================================
\begin{figure}[t]
\centering
% Training vs Inference Comparison
\begin{tikzpicture}[
    scale=0.9, every node/.style={scale=0.9}
]

% Training Phase (Left Side)
\begin{scope}
    \node[font=\small\bfseries] at (0, 3.5) {Training Phase};
    
    % Training boxes
    \node[draw, thick, fill=blue!15, rectangle, minimum width=35mm, minimum height=8mm, font=\footnotesize] (train_data) at (0, 2.5) {200 Scenarios (N=500)};
    
    \node[draw, thick, fill=orange!15, rectangle, minimum width=35mm, minimum height=8mm, font=\footnotesize, align=center] (blossom) at (0, 1.3) {Blossom Matching\\(Ground Truth)};
    
    \node[draw, thick, fill=green!15, rectangle, minimum width=35mm, minimum height=8mm, font=\footnotesize, align=center] (gnn_train) at (0, 0.1) {GNN Training\\60 Epochs};
    
    \node[draw, thick, fill=purple!15, rectangle, minimum width=35mm, minimum height=8mm, font=\footnotesize, align=center] (model) at (0, -1.1) {Trained Model\\(166K params)};
    
    % Arrows
    \draw[->, >=stealth, thick] (train_data) -- (blossom);
    \draw[->, >=stealth, thick] (blossom) -- (gnn_train);
    \draw[->, >=stealth, thick] (gnn_train) -- (model);
    
    % Time annotation
    \node[font=\tiny, gray, right=1mm of blossom] {30.4 s/graph};
    \node[font=\tiny, gray, right=1mm of gnn_train] {$\sim$2 hours total};
\end{scope}

% Arrow between phases
\draw[->, >=stealth, ultra thick, gray] (2, 0.5) -- (3.5, 0.5) node[midway, above, font=\small] {Deploy};

% Inference Phase (Right Side)
\begin{scope}[xshift=6cm]
    \node[font=\small\bfseries] at (0, 3.5) {Inference Phase};
    
    % Inference boxes
    \node[draw, thick, fill=blue!15, rectangle, minimum width=35mm, minimum height=8mm, font=\footnotesize] (new_csi) at (0, 2.5) {New CSI Snapshot};
    
    \node[draw, thick, fill=cyan!15, rectangle, minimum width=35mm, minimum height=8mm, font=\footnotesize, align=center] (fast_gnn) at (0, 1.3) {GNN Inference\\(No optimization)};
    
    \node[draw, thick, fill=green!15, rectangle, minimum width=35mm, minimum height=8mm, font=\footnotesize] (greedy) at (0, 0.1) {Greedy Matching};
    
    \node[draw, thick, fill=yellow!30, rectangle, minimum width=35mm, minimum height=8mm, font=\footnotesize, align=center] (result) at (0, -1.1) {User Pairs\\(96.5\% optimal)};
    
    % Arrows
    \draw[->, >=stealth, thick] (new_csi) -- (fast_gnn);
    \draw[->, >=stealth, thick] (fast_gnn) -- (greedy);
    \draw[->, >=stealth, thick] (greedy) -- (result);
    
    % Time annotation
    \node[font=\tiny, text=green!50!black, right=1mm of fast_gnn] {\textbf{752 ms}};
    \node[font=\tiny, text=green!50!black, right=1mm of result] {\textbf{40.4× faster}};
\end{scope}

% Bottom comparison box
\node[draw, double, thick, fill=gray!10, rectangle, text width=75mm, font=\footnotesize, align=center] at (3, -2.5) {
\textbf{Key Advantage:} Expensive optimization done once offline\\
Fast inference for real-time deployment
};

\end{tikzpicture}

\caption{Training versus inference phases. During training, expensive Blossom matching generates ground-truth labels for 200 scenarios. At inference, the trained GNN processes new CSI snapshots in 752 ms (40.4× speedup), demonstrating the value of learning-based optimization.}
\label{fig:training_inference}
\end{figure}

% ========================================
% REQUIRED PACKAGES
% Add to your preamble:
% ========================================
% \usepackage{tikz}
% \usetikzlibrary{shapes, arrows, positioning, calc, shapes.geometric}

% ========================================
% NOTES:
% ========================================
% 1. All figures use TikZ - ensure tikz package is loaded
% 2. Figures are designed for IEEE double-column format
% 3. Adjust scale parameter in each .tex file if needed
% 4. Colors are printer-friendly (tested in grayscale)
% 5. Font sizes are optimized for IEEE column width

% ========================================
% CUSTOMIZATION:
% ========================================
% - To change colors: modify fill=color!XX in each .tex file
% - To adjust size: change scale=X.XX parameter
% - To modify complexity: edit gray annotations
% - To change fonts: modify font=\footnotesize styles
